\documentclass[a5paper,12pt]{article}
\usepackage[utf8]{inputenc}
\usepackage[T1]{fontenc}
\usepackage[a5paper, top=2cm, bottom=2.5cm, left=2cm, right=2cm]{geometry}
\usepackage{titlesec}
\usepackage{parskip}
\usepackage{microtype}

\newcommand{\snapsname}[1]{\vspace{2em}{\footnotesize\scshape #1}\par\nopagebreak[4]\vspace{-1.5em}}
\newcommand{\melodi}[1]{\vspace{-\parskip}{\small\itshape Melodi: #1}\par\nopagebreak[4]\smallskip}

% Song section formatting
\titleformat{\section}[block]{\large\bfseries\centering}{}{0em}{}
\titlespacing{\section}{0pt}{1.5em}{0.5em}

\begin{document}
\centering

\begin{titlepage}
  \vspace*{4cm}
  {\Huge\bfseries Snapsvisor}\\[1em]
  \vfill
\end{titlepage}

\snapsname{Helan}
\section{Helan går}

Helan går,\\
sjung hopp falleri faderallan lej,\\
helan går,\\
sjung hopp faderallan lej.\\
Och den som inte helan tar,\\
han inte heller halvan får.\\
Helan går,\\
sjung hopp faderallan lej.

\snapsname{Halvan}
\section{Små nubbarna}
\melodi{Små grodorna}

Små nubbarna, små nubbarna er vill jag inte se.\\
Små nubbarna, små nubbarna er vill jag inte se.\\
I glasen på borden ska ni ej vara mer.\\
I glasen på borden ska ni ej vara mer.\\
För ni ska ner, för ni ska ner, för ni ska ner i mig,\\
För ni ska ner, för ni ska ner, för ni ska ner i mig.

\newpage

\snapsname{Tersen}
\section{Ingemar Bergman}

Tystnad. Tagning!



\snapsname{Kvarten}
\section{Jag hade en gång en snaps}
\melodi{Jag hade en gång en båt}

Jag hade en gång en snaps,\\
Jag hörde en ton slås an.\\
Så sjöng vi en liten sång\\
och snapsen försvann!\\
Svara mej du: Var är den nu?\\
Jag bara undrar: Var är den nu?


\snapsname{Kvinten}
\section{Tänk om jag hade lilla nubben}
\melodi{Hej tomtegubbar}

Tänk om jag hade lilla nubben uppå ett snöre i halsen,\\
tänk om jag hade lilla nubben uppå ett snöre i halsen.\\
Jag skulle dra den upp och ner,\\
så att den kändes som många fler.\\
Tänk om jag hade lilla nubben uppå ett snöre i halsen!

\newpage

\snapsname{Sexten}
\section{Sån är spriten}
\melodi{Sånt är livet}

Sån är spriten, ja sån är spriten,\\
Så mycket brännvin finns det här.\\
Den sup man ratar, den tar nån annan,\\
Så ta nu snapsen som du har kär.


\snapsname{Septen}
\section{Till spritbolaget ränner jag}
\melodi{Emils sång}

Till spritbolaget ränner jag\\
Och bankar på dess port.\\
Jag vill ha nåt' som bränner bra\\
Och gör mig sketfull fort.

Expediten fråga och sa:\\
Hur gammal kan min herre va?\\
Har du nåt legg ditt fula drägg,\\
Kom hit igen när du fått skägg.

Nej, detta var ju inte bra,\\
Jag ska bli full i kväll.\\
Då kom jag på en bra idé,\\
Dom har ju sprit på Shell.

Många flaskor stod där på rad.\\
Hurra, nu kan jag bli full och glad.\\
Den röda drycken rann ju ner.\\
Nu kan jag inte se nåt mer.


\snapsname{Rivan}
\section{Siffervisan}
\melodi{Ritsch, ratsch, filibom}

1 2 75 6 7 75 6 7 75 6 7\\
1 2 75 6 7 75 6 7 73\\
107 103 102 107 6 19 27\\
17 18 16 15 13 19 14 17\\
19 16 15 11 9 47 -- hej!


\snapsname{Räfflan}
\section{Idas snapsvisa}
\melodi{Idas sommarvisa}

Ni ska inte tro vi blir fulla\\
ifall inte vi sätter fart.\\
Och gör så att glasen blir fyllda\\
och flaskan så härlig och tom.\\
Sjunga och glädjas och kul ska vi ha\\
och festa till midsommardan.\\
Sen går vi hem med ett minne kvar\\
som varar till följande år.

\newpage

\snapsname{Rännan}
\section{Lille sup}
\melodi{Lille katt}

Lille sup, lille sup,\\
Lille söte supen.\\
Nu ska du, nu ska du,\\
Ner igenom strupen.

\snapsname{Smuttan}
\section{Hell and gore}
\melodi{Helan går}

Hell and gore,\\
Chung Hop father Allan Allan ley.\\
Hell and gore,\\
Chung Hop father Allan ley.\\
Oh handsome in the hell and tar,\\
hand hell are in the half and four.\\
Hell and gore,\\
Chung Hop father Allan ley.

\snapsname{Smuttans unge}
\snapsname{Lilla Manasse}
\snapsname{Lilla Manasses broder}
\snapsname{Femton droppar}
\snapsname{Kreaturens återuppståndelse}
\snapsname{Professorns klagan (Absolut sista supen)}
\snapsname{Den bleka dödens dryck}

\end{document}
